\documentclass{article}
\usepackage{amsthm,amssymb,mathtools,amsmath}

\newtheorem{theorem}{Theorem}[section]
\newtheorem{definition}[theorem]{Definition}
\newtheorem{notation}[theorem]{Notation}
\newtheorem{remark}[theorem]{Remark}

%\DeclareMathOperator*{\max}{max}
%\DeclareMathOperator*{\min}{min}

\begin{document}
\section{Cross Norms}
\begin{notation}
In this section, $A_1,$ $A_2$ will denote Banach spaces.
The norms on these spaces will be denoted by $\lVert \cdot \rVert$
unambiguously as we will distinguish between elements of $A_1$ and $A_2.$
\end{notation}
\begin{notation}[]
For Banach spaces $A_1, A_2,$ we will denote their \emph{algebraic tensor product} by $A_1 \otimes A_2.$
Its completion in some norm $\lVert \cdot \rVert_\beta$ will be denoted by $A_1 \otimes_\beta A_2.$
\end{notation}
\begin{definition}
	A norm $\lVert \cdot \rVert_c$ on $A_1 \otimes A_2$
is called a \emph{cross norm} if it satisfies $\lVert x_1 \otimes x_2 \rVert_c = \lVert x_1 \rVert \lVert x_2 \rVert$
for all $x_1 \in A_1, x_2 \in A_2.$
\end{definition}
\begin{remark}
	That $\lVert \cdot \rVert_\lambda$ is a cross norm is immediate from an application of the Hahn Banach Theorem.
\end{remark}
Suppose 
\begin{definition}
	The \emph{injective cross norm} $\lVert \cdot \rVert_\lambda$ on $A_1 \otimes A_2$ is defined by
	$$ \lvert |\xi|| = \sup \{ |\sum_{k=1}^{n} f(x_{1,k}) g(x_{2,k}) \rvert 
		\colon 
		\begin{array}{c}
		\xi = \sum_{k=1}^{n} x_{1,k}\otimes x_{2,k} = \xi, \\
		f \in A_1^\ast, \lVert f \rVert = 1, \\
		g \in A_2^\ast, \lVert g \rVert = 1
	\end{array}
\}.$$
\end{definition}

\begin{remark}
	Let $\lVert \cdot \rVert_\beta$ be a cross norm on $A_1 \otimes A_2.$ 
	Retaining the notation of the previous definition, 
	$$ \lvert f(x_1) \rvert \lvert g(x_2) \rvert
	\leq \lVert x_1 \rVert \lVert x_2 \rVert = \lVert x_1 \otimes x_2 \rVert_\beta.$$
\end{remark}
So, the injective cross norm is the smallest cross norm on $A_1 \otimes A_2.$

\begin{definition}
	The \emph{projective cross norm} on $A_1 \otimes A_2$ is denoted by 
	$\lVert \cdot \rVert_\gamma$ and defined as
	$$\lVert \xi \rVert = \inf \{ \sum_{k=1}^{n} \lVert x_{1,k} \rVert \lVert x_{2,k} \rVert
	\colon \xi = \sum_{k=1}^{n} x_{1,k} \otimes x_{2,k} \}.$$
\end{definition}

\begin{remark}
	Let $\lVert \cdot \rVert_\beta$ be a cross norm on $A_1 \otimes A_2.$ 
	For $\xi$ in $A_1 \otimes A_2,$ and $\epsilon>0,$ there exists a representation
	$\xi = \sum_{k=1}^{n} x_{1,k} \otimes x_{2,k}$ in $A_1 \otimes A_2,$ such that
	\[
		\lVert \xi \rVert_\beta \leq
		\sum_{k=1}^{n} \lVert x_{1,k} \otimes x_{2,k} \rVert_\beta
		= \sum_{k=1}^{n} \lVert x_{1,k} \rVert \lVert x_{2,k} \rVert
		\leq \lVert \xi \rVert_\gamma + \epsilon.
	\]
\end{remark}
So, the projective norm majorises all cross norms. 
That it is itself a cross norm is evident from the computation
$$ \lVert x_1 \rVert \lVert x_2 \rVert = 
\lVert x_1 \otimes x_2 \rVert_\gamma \leq \lVert x_1 \otimes x_2 \rVert_\lambda \leq \lVert x_1 \rVert \lVert x_2 \rVert.$$
The same computation reveals that all norms `between' the two norms are cross norms.
projective norm

\subsection{$C^\ast$ cross-norms}
In general, the projective cross\nobreakdash-norm is not a $C^\ast$~norm \cite[Remark 4.4]{Pisier_2020}.
So we consider 
\begin{definition}
	The maximal $\text{C}^\ast$ norm defined by
	\[
		\lVert x \rVert_{\max} = \sup \{
		\lVert \pi(x) \rVert \colon \pi \text{ is a representation of } A_1 \otimes A_2 \}
	\]

\end{definition}
The max norm is a cross norm.
\begin{remark}
	For $\xi = \sum_{k=1}^{n} x_{1,k} \otimes x_{2,k},$ in $A_1 \otimes A_2,$ we have
	$$ \lVert \pi(\xi) \rVert
	= \lVert \sum_{k=1}^{n} \pi(x_{1,k} \otimes x_{2,k}) \rVert
	\leq \sum_{k=1}^{n} \lVert \pi(x_{1,k} \otimes x_{2,k}) \rVert
	\leq \sum_{k=1}^{n} \lVert x_{1,k} \otimes x_{2,k} \rVert .$$
	This shows that $\lVert x \rVert_{\max} \leq \left\lVert x \right\rVert_\gamma.$

	In addition, for any representation $\pi_1$ of $A_1$ and $\pi_2$ of $A_2,$ we have 
	\begin{flalign*}
		\lVert \pi_1(x_1) \rVert \lVert \pi_2(x_2) \rVert 
		=\lVert \pi_1(x_1) \pi_2(x_2)\rVert
		= \lVert (\pi_1 \otimes \pi_2) (x_1 \otimes x_2) \rVert \\
		\leq \lVert x_1 \otimes x_2 \rVert_{\max} 
		\leq \lVert x_1 \otimes x_2 \rVert_\gamma
		\leq \lVert x_1 \rVert \lVert x_2 \rVert.
	\end{flalign*}
	By choosing faithful representations $\pi_1$ and $\pi_2,$ the left hand side and right hand side become equal.
\end{remark}

\begin{definition}
	The minimal $C^\ast$~cross\nobreakdash-norm $\lVert \cdot \rVert_{\min}$ is defined by
$$\lVert \xi \rVert_{\min} 
= \sup\{ \lVert (\pi_1 \otimes \pi_2) (\xi) \rVert \colon 
\begin{array}{c}
\pi_1 \text{ is a representation of } A_1, \\
\pi_2 \text{ is a representation of } A_2 
\end{array}
\} .$$
\end{definition}
It is immediate that $\lVert \cdot \rVert_{\min} \leq \lVert \cdot \rVert_{\max} \leq 
\lVert \cdot \rVert_\gamma.$ 
\begin{remark}
	To see that $\lVert \cdot \rVert_{\min}$ is a cross norm, we only need to show that $\lVert \cdot \rVert_{\Lambda}
	\leq \lVert \cdot \rVert_{\min}.$
	For $f$ in $A_1^\ast,$ there exist a representation $\pi_i$ of $A_1$ on a Hilbert Space $\mathcal{H}_1 $ and $\xi_1, \eta_1$
	in $\mathcal{H}_1$ satisfying $\lVert f \rVert = \lVert \xi_1 \rVert \lVert \eta_1 \rVert$ and
	$f(a) = \langle\eta_1, \pi_1(a) \xi_1\rangle$ for all $a$ in $A_1.$ Similarly, for $g$ in $A_2^\ast,$ there exist a representation
	$\pi_2$ of $A_2$ on a Hilbert space $\mathcal{H}_2$ and $\xi_2, \eta_2$ in $\mathcal{H}_2$ such that
	$g(b) = \langle \eta_2, \pi_2(b) \xi_2 \rangle$ for all $b$ in $A_2$ and $\lVert g \rVert = \lVert \xi_2 \rVert \lVert \eta_2 \rVert.$ This means for $x = \sum_{k=1}^{n} x_{1,k} \otimes x_{2,k}$ in $A_1 \otimes A_2,$ we have
	\begin{align*}
		{}& \lvert (f \otimes g) x \rvert \\
		={}&  \lvert \sum_{n=1}^{k} f(x_{1,k}) g(x_{2,k}) \rvert \\
		\leq{}&  \lvert \sum_{n=1}^{k} f(x_{1,k}) g(x_{2,k}) \rvert \\
		\leq{}&  \lvert \sum_{k=1}^{n} \langle\eta_1, \pi_1(x_{1,k}) \xi_1\rangle \rvert
		\lvert \langle\xi, \pi_2(x_{2,k}) \xi_2\rangle \rvert \\
		={}& \lvert \sum_{k=1}^{n} \langle\eta_1, \pi_1(x_{1,k}) \xi_1\rangle 
		\langle\xi, \pi_2(x_{2,k}) \xi_2\rangle \rvert \\
		={}& \lvert \sum_{n=1}^{k} \langle \eta_1 \otimes \eta_2, (\pi_1 \otimes \pi_2) (x_{1,k} \otimes x_{2,k}) (\xi_1 \otimes \xi_2) \rangle \rvert \\
		={}&  \lvert \langle \eta_1 \otimes \eta_2, (\pi_1 \otimes \pi_2) ( x) (\xi_1 \otimes \xi_2)\rangle \rvert \\
		\leq{}& \lVert \eta_1 \rVert \lVert \eta_2 \rVert \lVert (\pi_1 \otimes \pi_2)(x) \rVert
		\lVert \xi_1 \rVert \lVert \xi_2 \rVert \\
		={}& \lVert f \rVert \lVert g \rVert \lVert x \rVert_{\min}
	\end{align*}
	In other words, $\lVert x \rVert_\gamma \leq \lVert x \rVert_{\min}.$
\end{remark}
Thus, $\lVert \cdot \rVert_{\min}$ is a $C^\ast$ norm.


tensor products of ucp maps
functoriality
and universal property

\bibliographystyle{plain}
\bibliography{../notes/my_refs}

\end{document}

